একটি সেট~$A$-এর উপর \emph{সম্পর্ক}~$R$ হলো~$A$-এর
উপাদানগুলিকে পরস্পরের সঙ্গে সম্পর্কিত করার একটি উপায়। $x$ ও~$y$-এর মধ্যে
সম্পর্কটি \emph{থাকলে} আমরা $Rxy$ লিখি। আনুষ্ঠানিকভাবে, $R$-কে
সেই সব যুগল $\tuple{x,y} \in A^2$-এর সেট হিসেবে ধরা যায় যাদের
জন্য~$Rxy$ হয়। ছোট হওয়া, বড় হওয়া, সমান হওয়া, নিঃশেষে ভাগ করা,
সমান দৈর্ঘ্যের হওয়া, উপসেট হওয়া এবং সমান আকারের হওয়া—সবই
(সংখ্যা, প্রতীকক্রম বা সেটের সেটের উপর) সম্পর্কের গুরুত্বপূর্ণ
উদাহরণ। সম্পর্ককে চিত্রে দেখানোর একটি সাধারণ উপায় হলো \emph{গ্রাফ}।
আবার একটি গ্রাফকে একটি দ্বিপদ সম্পর্ক (\emph{প্রান্ত} সম্পর্ক) এবং
তার অন্তর্নিহিত \emph{শীর্ষগুলির} সেট হিসেবেও দেখা যায়।

কিছু সম্পর্কের এমন কয়েকটি ধর্ম আছে, যেগুলি তাদের বিশেষভাবে
আকর্ষণীয় বা উপযোগী করে তোলে। প্রত্যেক বস্তু নিজের সঙ্গে
$R$-সম্পর্কে থাকলে~$R$ \emph{প্রতিবিম্ব ধর্মবিশিষ্ট}। যেকোনো $x$ ও
$y$-এর জন্য $Rxy$ থেকে $Ryx$ অনুসৃত হলে সম্পর্কটি \emph{প্রতিসম}।
$Rxy$ ও $Ryz$ থেকে~$Rxz$ নিশ্চিত হলে সম্পর্কটি \emph{পরিযায়ী}।
যে সম্পর্কগুলির এই তিনটি ধর্মই আছে, সেগুলি \emph{তুল্যতা সম্পর্ক}।
$Rxy$ ও $Ryx$ থেকে~$x=y$ নিশ্চিত হলে সম্পর্কটি \emph{বিপ্রতিসম}।
প্রতিবিম্ব ধর্মবিশিষ্ট, বিপ্রতিসম ও পরিযায়ী সম্পর্কগুলিই
\emph{আংশিক ক্রম}। কোনো আংশিক ক্রমে যেকোনো $x$ ও~$y$-এর জন্য
$Rxy$ বা $x=y$ বা $Ryx$ হলে সেটি একটি \emph{রৈখিক ক্রম}।
(সাধারণভাবে, এই ধর্মবিশিষ্ট সম্পর্ককে \emph{সংযুক্ত} বলে।)

সম্পর্কগুলি (যুগলের) সেট, তাই সেটের ক্রিয়াগুলি সম্পর্কের উপরও
প্রয়োগ করা যায় (যেমন, সম্পর্কের সংযোগ ও ছেদ গঠন করা যায়)।
সম্পর্কগুলিকে পরপর জুড়েও দেওয়া যায় (\emph{আপেক্ষিক গুণফল}~$R \mid S$)।
$R$-এর সঙ্গে তার নিজের আপেক্ষিক গুণফল যতবার ইচ্ছা নিয়ে যে
যুগলগুলি পাওয়া যায়, সেগুলি নিয়ে~$R$-এর \emph{পরিযায়ী আবরণ}~$R^+$ গঠিত।
